\renewcommand{\thechapter}{\Asbuk{chapter}}

\chapter{Основные артефакты работы}

В этом приложении перечислены основные артефакты, полученные в ходе работы. Они включены сюда для воспроизводимости экспериментов и для того, чтобы отделить технические детали от основного текста.

\section{Репозитории и данные}

Исходный код проекта расположен в репозитории:
\begin{quote}
    \ttfamily https://github.com/halaction/humor-generation
\end{quote}

Подготовленный набор данных расположен в репозитории Hugging Face:
\begin{quote}
    \ttfamily https://huggingface.co/datasets/halaction/humor-generation
\end{quote}

Он содержит очищенный корпус шуток, векторные представления, ключевые слова, многореференсные запросы, кандидатные ответы, результаты попарной оценки и таблицы лидеров.

Обученные LoRA-адаптеры опубликованы отдельно от базовых моделей:
\begin{quote}
    \ttfamily https://huggingface.co/halaction/Qwen3-1.7B-humor-lora
\end{quote}
\begin{quote}
    \ttfamily https://huggingface.co/halaction/Qwen3-4B-humor-lora
\end{quote}

Такое разделение артефактов позволяет воспроизводить отдельные этапы: построение данных, обучение адаптера, генерацию кандидатов, автоматическую оценку и агрегацию таблиц лидеров.

\section{Структура программного конвейера}

Репозиторий организован как набор отдельных конвейеров и обучающих модулей. Основные компоненты приведены в Таблице~\ref{tab:appendix-code-structure}.

\begin{table}[h]
    \centering
    \begin{tabular}{p{5cm}p{9cm}}
        \hline
        Компонент                                      & Назначение                                                                                \\
        \hline
        \texttt{src/pipelines/jokes.py}                & Загрузка сырых корпусов, очистка, удаление дублей и сохранение корпуса шуток              \\
        \texttt{src/pipelines/embeddings.py}           & Построение плотных векторных представлений шуток                                          \\
        \texttt{src/pipelines/keywords.py}             & Извлечение ключевых слов с использованием сходства и MMR                                  \\
        \texttt{src/pipelines/references.py}           & Построение многореференсных запросов через векторный поиск                                \\
        \texttt{src/pipelines/candidates.py}           & Генерация кандидатных ответов разных моделей                                              \\
        \texttt{src/pipelines/evaluation.py}           & Попарная автоматическая оценка и построение таблиц лидеров                                \\
        \texttt{src/training/reference\_likelihood.py} & Вычисление правдоподобия эталонных ответов через принудительное учительское декодирование \\
        \texttt{src/training/mrvf\_trainer.py}         & Обучение LoRA-адаптера с MRVF-целью                                                       \\
        \texttt{scripts/train\_mrvf.py}                & Основной скрипт запуска обучения                                                          \\
        \hline
    \end{tabular}
    \caption{Основные компоненты программного конвейера}
    \label{tab:appendix-code-structure}
\end{table}

\chapter{Шаблоны запросов}

\section{Шаблон для построения эталонных запросов}

Для построения многореференсных запросов используется короткий шаблон на английском языке:
\begin{quote}
    \ttfamily Write a joke using the following keyword(s): \{keywords\}
\end{quote}

Этот шаблон используется как единая текстовая форма запроса. Он применяется при поиске эталонных шуток, генерации кандидатов и оценке ответов. Такое единообразие уменьшает несоответствие между построением набора данных и последующим сравнением моделей.

Главный недостаток шаблона состоит в его общей форме. Он не задаёт аудиторию, стиль, длину или конкретный юмористический механизм. Это сделано намеренно: работа изучает не узкий формат шутки, а способность модели использовать слабые ключевые слова и эталонные множества.

\section{Инструкции для векторного поиска}

Для кодирования ключевых слов и запросов в пространстве векторных представлений используются инструкции, ориентированные на поиск. Они помогают модели векторных представлений интерпретировать короткую строку не как изолированный текст, а как поисковый запрос.

Для ключевого слова используется шаблон:
\begin{quote}
    \ttfamily Instruct: Given a keyword for a joke, retrieve jokes that would match the keyword.\\
    Query: \{keyword\}
\end{quote}

Для полного запроса используется шаблон:
\begin{quote}
    \ttfamily Instruct: Given a joke request, retrieve jokes that would answer the request.\\
    Query: \{query\}
\end{quote}

Эти шаблоны не являются частью обучаемой языковой модели. Они используются на этапе построения данных, чтобы связать короткие ключевые слова и корпус шуток в общем семантическом пространстве.

\section{Шаблон генерации кандидатов}

После исправления протокола генерации модель должна возвращать только финальную шутку. В общем виде инструкция имеет следующий смысл:
\begin{quote}
    \ttfamily Write one short joke using the given keyword(s). Return only the final joke. Do not explain it.
\end{quote}

Это ограничение оказалось важным для качества оценки. Если модель возвращает объяснение, рассуждение или повторяющийся текст, модель-оценщик сравнивает уже не только юмористическую идею, но и нарушение формата ответа. Поэтому в финальных экспериментах оценивались только финальные шутки.

\section{Оценочный запрос}

При попарной оценке модель-оценщик получает исходный запрос и два кандидатных ответа. Системная инструкция просит выбрать ровно одного победителя и вернуть результат в формате JSON. На уровне смысла оценочный запрос имеет следующую структуру:
\begin{quote}
    \ttfamily You are comparing two jokes written for the same request. Choose the joke that is funnier for a broad audience. Avoid ties. Return strict JSON with the winner.
\end{quote}

Такой дизайн намеренно прост. Он не является полной человеческой рубрикой оценки юмора, но обеспечивает воспроизводимое попарное сравнение нескольких вариантов модели.

\chapter{Дополнительные сведения об экспериментах}

\section{Обучающие конфигурации}

Основные параметры финальных MRVF-запусков приведены в Таблице~\ref{tab:appendix-training-configs}. Эта таблица дублирует наиболее важные параметры из конфигурационных файлов и помогает воспроизвести экспериментальную постановку.

\begin{table}[h]
    \centering
    \begin{tabular}{lcc}
        \hline
        Параметр                      & Qwen3-1.7B         & Qwen3-4B           \\
        \hline
        Число шагов обучения          & 50                 & 70                 \\
        Базовая модель                & Qwen3-1.7B         & Qwen3-4B           \\
        Тип адаптера                  & LoRA               & LoRA               \\
        LoRA rank                     & 32                 & 16                 \\
        LoRA alpha                    & 64                 & 32                 \\
        Размер группы \(G\)           & 2                  & 2                  \\
        Число эталонов на запрос      & 2                  & 2                  \\
        Макс. длина рассуждения       & 512                & 384                \\
        Макс. длина эталонного ответа & 64                 & 64                 \\
        Тип целевой функции           & \(\log\)-масса     & \(\log\)-масса     \\
        Нормировка эталона            & среднее по токенам & среднее по токенам \\
        Нормировка преимущества       & GRPO               & GRPO               \\
        \hline
    \end{tabular}
    \caption{Параметры финальных MRVF-запусков}
    \label{tab:appendix-training-configs}
\end{table}

\section{Сравниваемые варианты моделей}

В автоматической оценке для каждой размерности модели сравнивались четыре варианта. Они приведены в Таблице~\ref{tab:appendix-model-variants}.

\begin{table}[h]
    \centering
    \begin{tabular}{lll}
        \hline
        Группа & Модель                          & Описание                       \\
        \hline
        1.7B   & \texttt{qwen3-17b-base-nothink} & Базовая модель без рассуждения \\
        1.7B   & \texttt{qwen3-17b-base-think}   & Базовая модель с рассуждением  \\
        1.7B   & \texttt{qwen3-17b-mrvf-nothink} & MRVF без рассуждения           \\
        1.7B   & \texttt{qwen3-17b-mrvf-think}   & MRVF с рассуждением            \\
        4B     & \texttt{qwen3-4b-base-nothink}  & Базовая модель без рассуждения \\
        4B     & \texttt{qwen3-4b-base-think}    & Базовая модель с рассуждением  \\
        4B     & \texttt{qwen3-4b-mrvf-nothink}  & MRVF без рассуждения           \\
        4B     & \texttt{qwen3-4b-mrvf-think}    & MRVF с рассуждением            \\
        \hline
    \end{tabular}
    \caption{Варианты моделей в финальной оценке}
    \label{tab:appendix-model-variants}
\end{table}
